\documentclass[12pt]{book}

\usepackage{graphicx}
\usepackage{eso-pic}
\usepackage{tcolorbox}
\usepackage{polyglossia}
\setdefaultlanguage[indentfirst=false,forceheadingpunctuation=false]{russian}
\setotherlanguages{english}

\RequirePackage[a4paper, headsep = 0.5 \headsep, left=2.5cm, right=2.1cm, top=2cm, bottom=2.1cm]{geometry}

\setmainfont{URW Bookman}
\newfontfamily\cyrillicfont{URW Bookman}[Script=Cyrillic, Ligatures=TeX]
\setlength{\parindent}{20pt}

\usepackage{enumitem}
\setlist[itemize, 1]{label =\raisebox{-0.3\height}{\scalebox{1.8}{\textbullet}}}

\title{Принимая отчужденность}
\author{Тодд МакГоуэн}
\date{2019}
\begin{document}

\AddToShipoutPictureBG*{\includegraphics[width=\paperwidth]{title.png}}

\null
    \thispagestyle{empty}%
    \addtocounter{page}{-1}%
    \newpage

\tableofcontents

\chapter{Введение}

''Отчужденность'' звучит непривлекательно. Будто что-то разрушительное, чего лучше избегать. Никто не говорит друзьям ''Давайте отчуждаться'', наверняка это не так приятно, как основание сообщества, романтические отношения, общение с приятелями, пицца, игра в теннис или выпивка. Напротив --- мы хотим помочь подростку избежать отчужденности, а отчужденному постороннему --- почувствовать себя уютней. Левые, согласно Карлу Марксу, понимают политическую деятельность как борьбу с отчуждением. Почти всем отчужденность видится проблемой, которую надо решить.

Приговор отчужденности поспешен и забывает о ее освободительном качестве. Исторически, у отчужденности плохая репутация, но ведь именно отчужденность дает возможность выхода за пределы конкретной ситуации. Следовало бы думать о ней гораздо лучше. Увидеть ее достоинства и постараться усилить ее опыт.

Никто открыто не отстаивал отчужденность как проект, но некоторые принимают ее основой свободы. Немецкий идеализм, кульминацией которого стала философия Гегеля, говорит об отчужденности как условии самосознания. Психоанализ воспринимает отчужденность субъекта как непреклонный факт нашего существования и стремится примирить с ним. Экзистенциализм Сартра и Симоны де Бовуар тоже видит в отчужденности субъекта основу свободы. Франц Фанон объединяет все три течения --- Гегеля, психоанализ и экзистенциализм --- и формулирует антиколониальную теорию, выступая против попыток избегать фундаментальной нам всем отчужденности. Впрочем, эти мыслители никогда не используют отчужденность как основное понятие или лозунг, хотя оно объединяет эти разрозненные попытки говорить о свободе. В этой книге отчужденность выдвигается на первый план.

Отчужденность парадоксальна. Она не значит отклонения от исходной идентичности, от того, что мы когда-то имели но потеряли. Отчужденность неизменно присуща нам. Она не отдаленность от первичного, потому что первична. Мы отчуждены в своей субъектности до самой смерти. Смерть --- единственное лекарство от этой болезни, так что, наверное, не такая уж это и болезнь. Отчужденность --- не то, о чем можно сожалеть или чего лучше избегать, ведь благодаря ней возможно то, что делает существование достойным продолжения --- свобода, равенство, солидарность, способность выходить за пределы ситуации.

Все творение зависит от отчужденности, она выкорчевывает нас из определенности, освобождает от конкретной ситуации. Каждый рождается в специфичном социуме, в котором на него действуют мощные силы, но не сводится к ним. Первичное отчуждение в языке дистанцирует как от других, так и внутренне: к себе теперь можно относиться как к другому. Мы отчуждены от себя и от других, и что бы мы ни делали, нам никогда не преодолеть эту отчужденность. Она основа нашего существования.

Отчужденность --- это отсутствие самоидентичности. Она отдаляет субъекта от условий, в которых он проявляется. Если я не тождествен себе, если я в противоречии с собой, то внешние силы не могут меня полностью определить. Существо, дистанцированное от самого себя, никогда не сможет быть тем, во что его толкают внешние силы. Этот внутренний раскол значит сопротивление детерминированию биологией или обществом. Именно дистанция от самой себя позволяет субъектности действовать против определяющего. Социальное и натуралистическое утрачивают свою власть на субъектом, потому что его отчужденность делает их незначительными.

Я могу не оставаться до конца своих дней жителем маленького американского городка на Среднем Западе, могу бросить вызов определяющим факторам этой ситуации, и все потому что никогда не отождествлен с положением в котором родился или нахожусь. Это дает мне возможность отнестись к ситуации критически, вплоть до отказа от всех идей и ценностей, которые возложило на меня общество. Я могу быть воспитан с энтузиазмом к капиталистическому строю, но стать марксистом. Если меня заставляют быть женственным, я могу стать мужественным. Полюбив Средний Запад Америки, я, возможно, стану работать в Аргентине. Моя неспособность быть именно тем, кто я есть --- моя не-однозначность --- вот что позволяет мне выйти из исходной ситуации. И даже если я остаюсь в ней и принимаю ее, я все равно отношусь к ней как отчужденный. Можно быть чужим и у себя дома. У меня, как и у всех, есть внутренняя чужеродность, мешающая прямому усвоению диктата общества.

Без этой внутренней дистанции человек не был бы субъектом, не смог бы соотноситься с собой. Субъектность порождена и определена этим расколом. Отделенность от самих себя делает нас нами --- неотделимыми от чуждого в нас. Чуждое дает нам нашу уникальность.

Эта внутренняя чуждость и есть то, что Зигмунд Фрейд назовет бессознательным. Бессознательное ставит барьер как между субъектом и ситуацией, так и внутри субъекта. Только отчужденные обладают бессознательным. Вернее, чуждый статус бессознательного значит, что это оно владеет нами. Это иное внутри диктует субъекту его действия, направляет усилия и создает убеждения. Бессознательное ориентирует субъектность, потому что выражает ее самораскол.

Отчуждение субъектов от себя видно, когда бессознательное заставляет их идти наперекор своим же интересам. Кто-то объедается пончиками и кексами, вместо того, чтобы есть яблоки и брокколи и предотвратить ишемическую болезнь. Другой вместо комфортного, неконфликтного партнера выбирает сложного, болезненного человека, с которым постоянно ссорится. Продолжает искать и употреблять все новые дурманящие вещества, хотя прекрасно понимает, как они опасны и вредны. Хочет профессионального роста, но раз за разом выступает против начальника. Повсюду субъекта поджидают саморазрушительные, ускользающие от сознательного контроля действия.

Каждый сосуществует с идущей вразрез его желаниям силой. Каждый видит, как сам мешает достижению своих целей, как не может контролировать себя. Отсутствие контроля изнурительно, но также основа способности субъекта становиться иным. Верховенство чуждого --- бессознательного --- открывает возможность свободы. Свобода напрямую зависит от подрыва сознательных желаний бессознательными импульсами.

О бессознательном обычно думают как о принуждающем, забирающем свободу. Мы не обдумываем бессознательные действия и не чувствуем, что они выбраны нами. Но свобода --- не плод сознательных размышлений. Это прерывание решимости, определенности. И именно это делает бессознательное. Оно перебивает социальные и психологические определяющих факторы, открывает разрыв в необусловленную ими альтернативную возможность. Избегая сознательного контроля, бессознательное становится точкой свободы субъекта. Мы свободны благодаря своей раздвоенности.

Самотождественности лишена не только субъективность. Ничто не самотождественно, иначе не смогло бы меняться во времени и перемещаться в пространстве. Это касается и животных и растений, и гор и рек. Будь сущность тем, что она есть, она бы никогда не изменилась. Субъект же превращает свою нетождественность в определяющую его структуру, такой процесс и есть отчуждение.

Первое свидетельство отчуждения --- имя. Ни один ребенок не дает его себе сам. Происходит акт насилия, родители или другие взрослые навязывают ребенку произвольное имя. И даже если он потом воспринимает его своим, оно остается чуждым, сошедшим на него извне. Родители могут называть мальчика Эрнестом и надеется, что он вырастет честным и прямодушным. Могут называть девочку Грейс и думать, что она будет приличной и изящной. Иногда используют имена художников, таких как Лэнгстон или Дэшил, пытаясь ассоциацией приобщить детей к их гениальности. Дают имя, связанное с семейной историей, чтобы оно напоминало ребенку о его происхождении, определяло его поведение. Но и ничего не значащее имя остается чуждым названному означающим. Он не выбирал его, и воспринимает его своим только через отчуждение.

Имя это наша дистанция до себя. В каждом языке есть способ идентифицировать себя именем: я могу сказать ''Меня зовут Тодд'', или ''Ich heiße Todd'', или ''Me llamo Todd''. Всякое такое высказывание провозглашает тождество, и этим выявляет внутреннюю дистанцию между говорящим и именем. Так оно выдает, что я и Тодд --- не одно и то же. Самотождественное существо не смогло бы так сказать.

Даже если я переименую себя, имя останется чуждым. Казалось бы, я называю себя сам, не полагаюсь на данное другими. Разве этого недостаточно для преодоления отчужденности? Нет. Называя себя иначе, я действую исходя их доступных возможностей: выбираю из используемых другими имен или комбинирую их. Выбор имени всегда предполагает долг перед внешним. Мое новое имя не только принадлежит другим, но и выбрано с учетом их впечатления. Откуда его не возьми, оно несет печать самораскола.

Неизбежность отчужденности дает некоторую уверенность. И вместо того, чтобы страдать от саморазделения до самой смерти, мы можем манипулировать им и даже усиливать его. Оставление тщетных попыток преодолеть отчужденность, ее принятие, трансформирует ее. Из бремени она становится освобождением от определенности. Переосмысление отчужденности --- важнейшая политическая задача, особенно в эпоху вездесущих схем его якобы преодоления.

Об отчужденности плохо думают, потому что воображают скрывающуюся за ней завершенность. Где-то в затерянном прошлом, в Эдеме, или в утопическом будущем, нас ждет самотождественность, полнота, совершенство. Она всегда ускользает и тем соблазняет еще больше. Такое лишенное всякого отчуждения состояние очень привлекает нас, но лишь пока остается воображаемым. Малейшая попытка его очертить, конкретизировать, воплотить, выдает его полную несостоятельность.

Нам следует мыслить отчужденность вне ассоциаций с полнотой или самотождественностью. Мы отчуждены, но не от чего-то могущего дополнить нас. Мы и не теряли ничего такого. Помыслив нехватку вне связи с целостностью, мы поменяем политическую значимость отчужденности. Отчужденность не определяется через не-отчужденность.

Отчужденность --- источник всех человеческих страданий, но оно и источник всякого удовлетворения. Отчужденность предоставляет пространство действию вопреки, без него субъект неотличим от сил, породивших его. Заниматься сексом, есть торт или обсуждать немецкий идеализм --- ничего из этого не возможно без отчужденности, потому что без нее нет субъекта.

Если что и отличает говорящих существ от прочих, так это их осознание отсутствия самоидентичности. Через язык отчужденность осознается субъектами неизменно присущей им. Язык создает дистанцию, и это меняет отношение существа к отчужденности. Субъекты означающего могут осознать и принять себя отчужденными, могут найти удовлетворение в этом отсутствии самотождественности, и только отказ рассматривать отчужденность конститутивной составляющей мешает этому.

Чаще мы думаем о своих экзистенциальных и политических проектах как о преодолевающих отчужденность. За редкими исключениями, мы стремимся дать отчужденности бой, искоренить ее. Терапевты воображают, что помогают пациентам обрести самоидентификацию. Революционеры стремятся построить общество, в котором жили бы в гармонии с другими. Экоактивисты работают на формой существования, которая была бы в гармонии с природой. Разительное разнообразие проектов не затрагивает их сути --- избавления от чуждого. Отчужденность непрестанно видится тем, с чем следует бороться, хотя сама борьба возможна только благодаря ней.

Эта книга предлагает совершенно иной взгляд на отчужденность: усилия в ее преодолении значат не радикальный ответ ситуации, а неспособность увидеть конституирующую роль отчужденности. Нам следует оставить попытки преодолеть ее. Отчужденность должна стать нашей экзистенциальной и политической программой. Отчуждение --- это освобождение.

\end{document}
